Let $K=\left(\dfrac{1}{2}\cdot\dfrac{3}{2}\cdot\dfrac{5}{2}\cdot\dfrac{7}{2}\cdot\dfrac{9}{2}\right)^2$. At $x=\dfrac{11}{2}$ the product equals $-K$, so this is a root. Pairing factors symmetrically about this point shows the polynomial is a polynomial in $\left(x-\dfrac{11}{2}\right)^2$; consequently this central root has multiplicity at least two.

In each interval $(1,2),(3,4),(7,8),(9,10)$ the product is negative and zero at the endpoints. At the midpoint of $(3,4)$, its absolute value equals
\[
\left(\dfrac{1}{2}\dfrac{3}{2}\dfrac{5}{2}\right)^2
\dfrac{7}{2}\dfrac{9}{2}\dfrac{11}{2}\dfrac{13}{2}>K.
\]
At the midpoint of $(1,2)$ it is still larger, since its ordered positive distance factors dominate those at $3.5$. Reflection about $5.5$ gives the same statements in the other two intervals. The intermediate value theorem therefore gives at least two roots in each of these four intervals, eight in total.

Together with the central root counted at least twice, this exhausts the degree ten. Thus there are exactly $\boxed9$ distinct real roots.
